\section{Introduction}

IRES elements recruit translation machinery independently of the canonical cap-dependent pathway
and are attractive components for synthetic and therapeutic RNA systems. Their activity depends on
more than a short sequence motif: alternative RNA structures, host factors, construct design, cell
type, and downstream coding context can all alter translation.

Recent IRES-AI work established useful same-task prediction, evolutionary mutation, and diffusion
generation baselines using 46,774 labelled sequences. However, 95.4\% of that identification set
comes from a 174-nt DNA/lentiviral reporter library. RNA-based IRES-TrAPPr measurements and recent
IRES--cargo studies motivate an explicit audit of whether this legacy signal transfers to direct-RNA
and full-length contexts.

We therefore formulate IRES design as a constrained optimization problem rather than a pure
classifier-maximization task: can a candidate retain the energy and ensemble-structural features
of an experimentally supported full-length parent while improving a released IRES-recognition
score? Our contributions are: (i) \textsc{StructIRES-Rank}, an assay-aware, generator-agnostic
rank-selection layer that combines a sequence score, parent-relative energy preservation and
ensemble-structure preservation; (ii) a controlled full-length design matrix in which released
RNA-FM score-only, energy-only, ensemble-only and combined selection operate on exactly the same
candidates and budget; (iii) an independent direct-RNA held-out proxy that tests whether the
combined constraints improve computational candidate enrichment beyond each score-only or
single-constraint alternative; and (iv) a cross-assay benchmark audit that explains why legacy
IRES classifiers are useful baselines but insufficient standalone design objectives.
